% r03_h1s_b: proof of the two-flow comparison bound, part 2: the sets of non-positive bare degree, the normal form, the subtraction rule.

\begin{lemma}[the connected all-$C$ sets of non-positive bare degree]\label{rot:h1s:CL}
Let $H$ be a connected set of all-$C$ lines with every current-vertex derivative placed on an internal line of its unit when the vertex belongs to a unit (the placement (NF2) below), and let $n_e$, $n_s$, $n_\lambda$ be the numbers of its current, seagull and quartic vertices, $n_{\rm ct}$ the number of its counterterm insertions, $E$ the number of its slots (legs of $V(H)$ not on lines of $H$, external fields included) and $x$ the number of current-vertex derivatives of $V(H)$ placed outside $H$. Then
\[
D_0(H)=\tfrac12n_e+n_s+n_\lambda+n_{\rm ct}+\tfrac12E-3+x .
\]
Decompose $H$ into its two-edge-connected blocks $P$ joined by $\rho-1$ bridges $b$. Then $D_0(H)=\sum_P[\omega(P)+\mathrm{out}(P)]+\sum_b(2-\mathrm{in}(b))\ge\rho-2$, where $\omega(P)$ is the degree of $P$ with every derivative of its vertices placed inside $P$, $\mathrm{out}(P)$ the number of derivatives of vertices of $P$ placed outside $P$ and $\mathrm{in}(b)\le2$ the number of derivatives on the bridge, so that $\sum_P\mathrm{out}(P)\ge\sum_b\mathrm{in}(b)$; and $\omega(P)\ge-1$ for every block with legs, with equality exactly for the units SE and VP (the tadpoles being dropped), for a counterterm insertion and, by convention, for a bare current vertex ($\omega=-1$, $\mathrm{out}=1$). Consequently $D_0(H)\le0$ forces one of: $\rho=1$ and $H$ is a unit ($-1$) or one of T, SC, W, the double seagull, a $\lambda^2$ or $e^4$ or $e^2\lambda$ two-point structure ($0$); $\rho=2$ and $H$ is two blocks with $\omega=-1$ and every non-internal derivative on the bridge: $U$--$l$--$U'$ ($l$ undifferentiated), $U$--$\partial l$--$v$, $v$--$\partial\partial l$--$v'$ ($=$ L24), and with a counterterm insertion ct--$l$--$U$, ct--$\partial l$--$v$, ct--$l$--ct; and, in the local-term monomials of the rule below, the same with a vertex of weight $-1$ (the vacuum-polarisation leftover of \Cref{rot:h1s:VP}) joined by an undifferentiated photon line to a vacuum-polarisation compound whose leg is that line, or to another such vertex. Unions of blocks at a cut vertex (a quartic or seagull vertex) have both slots at one vertex and are common numbers. Every member is a two-point structure except T and SC (three slots) and the one-point numbers.
\end{lemma}
\begin{proof}
Count: $\sum_ld_l=-I-n_\partial$ with $I$ the number of lines and $n_\partial$ the number of internal derivatives, $n_\partial=n_e-x$; the scalar legs give $2I_\phi+E_\phi=2n_e+2n_s+4n_\lambda+2n_{\rm ct}$ and the photon legs $2I_A+E_A=n_e+2n_s$; $|V|=n_e+n_s+n_\lambda+n_{\rm ct}$, and the counterterm insertions carry the vertex weights $\sum_va_v=-n_{\rm ct}$. Substituting in $-I-n_\partial+\sum_va_v+3(|V|-1)$ gives the formula. Blocks and bridges: $3(|V(H)|-1)=\sum_P3(|V(P)|-1)+3(\rho-1)$ because the block--bridge graph is a tree; a scalar bridge has $d_b=-1-\mathrm{in}(b)$ and a photon bridge carries no derivative in the placement (NF2); a derivative of a vertex of $P$ placed on a bridge counts $+1$ in $\mathrm{out}(P)$ and $-1$ on the bridge, and every derivative on a bridge is the own derivative of one of its end vertices, each of which lies in a block, whence $\sum_P\mathrm{out}(P)\ge\sum_b\mathrm{in}(b)$ and $D_0(H)\ge-\rho+2(\rho-1)=\rho-2$. The formula applied to a block with all derivatives internal gives $\omega(P)=\frac12n_e+n_s+n_\lambda+\frac12E-3$ for a block without insertions, and a counterterm insertion alone is a block with $E=2$ and $\omega=-1$; hence $\omega(P)\ge-1$ with the stated equality cases ($E=2$: $(n_e,n_s,n_\lambda)=(2,0,0)$; $E=3$: a single current vertex, whose derivative is necessarily outside), and the enumeration of the pairs $(P,P')$ with $\omega=-1$ and $\mathrm{in}(b)\le2$ gives the $\rho=2$ list; $\rho\ge3$ is impossible. The connected sets of bare degree $0$ with all derivatives internal are those with $\frac12n_e+n_s+n_\lambda+\frac12E=3$: $E=2$ with $(4,0,0)$, $(0,0,2)$, $(2,0,1)$, $(0,2,0)$, $(2,1,0)$, $(0,1,1)$; $E=3$ with $(3,0,0)$, $(1,1,0)$, $(1,0,1)$; $E=4$ with $(2,0,0)$, which is L24 (two current vertices joined by one scalar line carrying both derivatives; joined by a photon line both derivatives are outside and the degree is $2$). Integrality: $E\equiv n_e\pmod2$.
\end{proof}

\begin{corollary}[placement at the terminals; the derivative count]\label{rot:h1s:nd}
\textup{(i)} In a connected set with $D_0(H)\le0$ no vertex of $H$ places its own derivative outside $H$. \textup{(ii)} After the rule of \Cref{rot:h1s:rule} every line carries at most one lattice derivative at each end (its end vertex's own derivative or one transferred derivative, never both), hence at most two in all. The leg of a compound carries at most one as formed; when the compound is the differenced factor of (R2b) at one of its outer vertices, its leg carries the difference (\Cref{rot:h1s:U}(iv)) and then at most two; the middle line of a double compound carries none; a photon line carries only transferred derivatives, at most one per end; a contour line receives at most one, at its bond end.
\end{corollary}
\begin{proof}
(i) With $x\ge1$ the formula gives $D_0(H)\ge-1+x\ge0$, with equality only if the set with internal placement is a unit; but a unit's vertices have their derivatives inside by (NF2). (ii) Own derivatives: a current vertex has one, placed on one of its two scalar lines. Transferred derivatives land on a line $l$ at a vertex $w$ of the component $H$ being processed, $l\notin H$ later than $H$, $D(H)=0$ at formation, hence $D_0(H)\le0$ and, by (i), $w$'s own derivative is inside $H$: $l$ carries no own derivative of $w$. A second transfer onto $l$ at $w$ would come from a component $H'\ni w$ formed before $l$ enters the segment; $H'\supseteq H$ (nested initial segments), so $D(H')\ge D_0(H')+1\ge1$ unless $H'$ is a bare unit, which is never a component (\Cref{rot:h1s:Dp}(b)); so $H'$ is not processed. The leg of a compound is attached to the unit at a vertex whose own derivative is internal and at which no rule acts (the unit's lines are not elements before the compound enters, and afterwards they are inside it), so the leg carries at most one derivative, at its far end; the middle line of a double compound has two such ends. Photon lines carry no own derivative (\Cref{rot:h1s:cv}). For contour lines see (iii) of the rule.
\end{proof}

\begin{definition}[normal form and the subtraction rule]\label{rot:h1s:rule}
In each flow, in its own frame, before any element is mixed:
\begin{enumerate}[label=\textup{(NF\arabic*)},leftmargin=*]
\item Tadpoles and their counterterm shares cancel exactly per slice (TAD): they are dropped.
\item At every vertex of a unit (SE, VP, T, SC) whose derivative sits on a leg, the derivative is written on the internal scalar line by \Cref{rot:h1s:cv}; the remainder term, in which the derivative acts on the vertex's photon field, is the unit's own photon line, so it stays inside the unit (for SE the kernel is the sum of the pieces $\partial\partial C\cdot D$, $\partial C\cdot\partial D$, $C\cdot\partial\partial D$, each with $\sum_ld_l=-4$). After this a unit's leg carries at most the far vertex's derivative.
\item The seagull and derivative tadpoles of the order-$e^2$ vacuum polarisation are eliminated through the exact Ward sum of \Cref{rot:h1s:VP}; what remains at a vacuum-polarisation unit is the bubble for every slice assignment plus one common vertex weight.
\item The telescoping is done per ordering class of \cite[(2.7)]{Ba3}. The two classes that differ only by swapping the two internal lines of a vacuum-polarisation unit are merged before telescoping (identical element sets; \Cref{rot:h1s:MS} is applied per class in absolute value, a factor $2$ per such unit), so a vacuum-polarisation slice assignment is the unordered pair $\{j_1,j_2\}$ and its kernel the line-symmetrised one. The two Wick contractions of a self-energy between the same two legs lie in one class and are one element (their sum is the symmetrised kernel of \Cref{rot:h1s:M}).
\end{enumerate}
Fix an ordering of the scale-summed elements and process the initial segments $G_1\subset G_2\subset\dots$. Let $\sigma_{\max}(U)$ be the coarsest scale of a unit $U$.
\begin{enumerate}[label=\textup{(R\arabic*)},leftmargin=*]
\item\emph{Compounds.} A unit $U$ (SE, VP; order two) or a degree-zero one-loop piece (T, SC; order one) one of whose vertices $y$ has all its non-$U$ elements later than $\sigma_{\max}(U)$ (an external field counts as later) is bundled into the \emph{compound} $(U,F_y)$: the own-frame Taylor subtraction, about $U$'s other vertex $x$ (in leg-site coordinates, the vertex sites shifted by the bonds of the currents), of the product $F_y$ of the elements at $y$, taken on the earliest line $f$ of $F_y$ (or on the external if $F_y$ has no line), to second order for SE, VP and to first order for T, SC; one leg per unit, preferring a derivative-free, unclaimed leg. In the reduced ordering the compound is one element at the position of $f$ with the vertex set $V(U)\cup V(f)$ and the line dimensions of $U$ and $f$: $D(\{(U,f)\})=2+d_f\in\{1,0\}$ (SE, VP), $=3+d_f\ge1$ (T, SC); with an external leg $D=\alpha$ resp. $1$ (\Cref{rot:h1s:U}). Two units forced onto the same leg $l$ form the \emph{double compound} $(U_1,l,U_2)$, one element at $l$'s position with $D=D_0(U_1\cup l\cup U_2)=0$. The local terms of the subtraction are exact and taken in the flow's own frame: for SE the zeroth-order term $c_\sigma f(x)g(x)$ cancels the co-switched counterterm share of the same $\sigma$ (\cite[(3.33)]{K1} through order four; \Cref{rot:lem:ctsame}) and the first-order term vanishes by \Cref{rot:h1s:M}; for VP the zeroth-order term enters the Ward sum of \Cref{rot:h1s:VP} and the first-order term vanishes by \Cref{rot:h1s:M}; for T the local term is $f(x)\sum_zT''(x,z)g(z)$ with the two-point kernel $T''(x,z)\defeq\sum_yT(x;y,z)$ of dimension $-3$, a degree-zero line processed by (R2); for SC the local term vanishes or carries an explicit factor $\eta'$ (the bond shift). A unit none of whose vertices qualifies (both legs earlier than $\sigma_{\max}$, the \emph{split} unit) is not an element: its lines are ordinary elements and its counterterm share stands alone as a loop line of dimension $-1$ at scale $\sigma_{\max}$ in a separate monomial.
\item\emph{Components of degree zero.} Whenever a newly formed connected component $H$ of the current initial segment has $D(H)=0$ (degrees after the earlier transfers), apply the first available clause.
\begin{enumerate}[label=\textup{(\alph*)},leftmargin=*]
\item If $H$ contains a bare current vertex $w$, or a vertex $w$ of a split unit, whose derivative sits on a line of $H$, whose other scalar field is not in $H$ and whose photon field is an external field or a photon line not in $H$: rewrite the vertex by \Cref{rot:h1s:cv}. This gives two monomials, (i) with the derivative on the other scalar field and (ii) with $\partial^*$ on the photon field; in both $D(H)$ rises to $1$, the recipient (later than $H$, since $H$ is a component) changing $d\mapsto d-1$, or at no cost if it is an external field; the term with $\partial^{\eta'}g$ carries an explicit $\eta'$, a vertex weight $1$. Each vertex is rewritten at most once. If $w$'s photon field is a line inside $H$, the clause is not available (the move would leave $D(H)=0$): this happens exactly for the outer vertices of the two-point structures W and for the leg-carrying vertices of a triangle, which go to (b) and (R1).
\item Else, if $H$ is a two-point structure (exactly two vertices $x$, $t$ of $H$ carry elements outside $H$ or externals): the first-order subtraction at the terminal $t$ about $x$ (either terminal may be chosen),
\[
\textstyle\sum_tH(x,t)G(t)=c\,G(x)+\sum_tH(x,t)[G(t)-G(x)],\qquad c\defeq\sum_tH(x,t),
\]
$G$ the product of the later elements at $t$, Leibniz-expanded so that exactly one factor is differenced. $c$ is a common number of both flows (\Cref{rot:h1s:Z0}), bounded by \Cref{rot:h1s:MS} applied per flow to $H$ with one vertex fixed ($|\sum_\sigma c_\sigma|\le C(1+r)^{n_0+1}$); for the sunset and the scalar $e^4$ structures it is $0$ exactly with the co-switched order-four counterterm share, for the $e^4$ vacuum polarisation it is $c\,\delta_{\mu\nu}$, for a triangle kernel $T''$ it is $0$ exactly (the summand is odd under the lattice inversion through a site: the current on a bond goes to minus the current on the inverted bond, the propagators are even and the three lines are distinguishable), and for a double compound standing alone it is $0$ exactly (translation invariance of the outer unit's kernel). $D(H)$ rises to $1$ (the factor $|t-x|$ absorbed on a line of $H$); the differenced factor changes $d\mapsto d-1$ if it is a line (for a compound, its leg, \Cref{rot:h1s:U}(iv)), at no cost if it is an external, a coarse element or a vertex function; the $Y$-displacement $|\tau t-\tau x-(t-x)|\le2c_\tau\eta'$ against the $Y$-gradient is an explicit-$\eta'$ vertex weight, a degree-zero component with gain $\eta'\ge\eta'^\theta$, counted among the logarithms.
\item Else $H$ has at most one vertex with outside elements or externals and is a common number (\Cref{rot:h1s:Z0} with $a=0$; both flows evaluate the same closed lattice sum), bounded per flow by \Cref{rot:h1s:MS}: a vertex weight $0$, never a switched element.
\end{enumerate}
\item\emph{Local-term monomials.} The monomials produced by the zeroth-order terms of (R1) and (R2b) (the unit or structure collapsed to a vertex of weight $0$, or $-1$ for the vacuum-polarisation leftover of \Cref{rot:h1s:VP}) are processed by the same rules; their degree-zero components are L24-type lines, the weight classes of \Cref{rot:h1s:CL}, and, when a triangle inside a W is collapsed, the kernels $\{\partial C,D\}$ and $\{C,\partial C\}$ between the collapsed point and the other terminal, all handled by (R2b) with common numbers.
\end{enumerate}
\textup{(iii)} At $e>0$ the recipient of (R2a)(ii) may be a contour line \cite[(3.27)]{K1} whose bond end sits at $w$; it receives this one derivative at its bond end and no other, since its contour end is the site of a contour vertex, which carries the weight $\ge\frac12$ and so never lies in a degree-zero component. Contour vertices are never recipients of any clause, for the same reason, and the constant-background phases \cite[(3.46)]{K1} are not photon fields. The lifted external background at $w$ is an external field with the gradient bound $\le Cp_r$ (\Cref{rot:lem:CB}(0)(R)).
\end{definition}

The rule is Ba{\l}aban's own device for the graph (2.4), ``the first few $G_1,G_2^{(2)},\dots$ are graphs (2.4)\dots replaced by graphs with degree $+1$'' by the integration by parts \cite[(2.8)]{Ba3}, applied to every degree-zero component that forms in an initial segment, with composite lines (compounds) allowed and with the first-order subtraction (R2b) for the components without a usable bare vertex. Two features are forced by the two flows: the units are bundled with a leg as one element, so that no line of a unit is ever processed alone inside an isolated unit (processing the lone $\partial\partial C$ of a self-energy with both legs later would lose the gain, since the Taylor order at the slot would be fixed before the second unit line enters), and the local terms are common numbers of both flows, so that no first moment of a flow-dependent object is ever mixed between the frames.
